\ssr{РЕФЕРАТ}

Расчетно-пояснительная записка 41 с., 18 рис., 4 табл., 9 источн., 1 прил.

СИСТЕМЫ ЧАСТИЦ, ВИЗУАЛИЗАЦИЯ ОГНЯ, ВНЕШНЕЕ ОСВЕЩЕНИЕ, МОДЕЛЬ ОСВЕЩЕНИЯ, Z-БУФЕР, ПРОИЗВОДИТЕЛЬНОСТЬ.

Цель работы – разработка программного обеспечения для визуализации трёхмерного огня методом систем частиц с учётом параметров внешнего освещения и исследование влияния количества частиц и параметров системы на производительность.

Для моделирования и отображения пламени использовался метод систем частиц, эмиссионная модель освещения с учётом внешнего точечного источника света и алгоритм z-буфера с аддитивным смешиванием для решения задачи видимости.

Разработаны формальные модели данных и схемы ключевых алгоритмов: основного цикла, обновления системы частиц, вычисления цвета частицы с учётом освещения, визуализации сцены. Реализовано спроектированное программное обеспечение на C++ с использованием фреймворка Qt, обеспечивающее интерактивное управление параметрами систем частиц, освещения и камеры в реальном времени. Описана структура сцены и форматы данных.

Проведено исследование производительности: проанализирована зависимость частоты кадров от общего количества частиц в сцене (от 10 000 до 2 000 000) и оценено влияние на производительность использования внешнего освещения и алгоритма z-буфера.

\clearpage
